% GENERATED FILE. Edit canonical lessons, then run npm run generate:books.
% canonical-source-hash: fnv1a64:39d523d4e4176588
% canonical-lessons: ML-C13-shareera-bhaagangal, ML-S08-letter-ta, ML-S117-letter-tta, ML-C13-checkpoint

\chapter{Head and Hand}
\label{ch:ml-head-and-hand}
\hlchaptermodality{13}

\begin{chapteropening}
\textbf{By the end of this chapter:} \emph{I can name the head and the hand in Malayalam.}

Complete ML-C13-checkpoint, the chapter's closing checkpoint: say tala and kai and match them with Tamil's talai and kai, find \ml{ത} in \ml{തല} and \ml{ട} in \ml{രണ്ട്}, then write both letters from sound.
\end{chapteropening}

\section[\texorpdfstring{\ml{തല} \ml{കൈ} (tala kai)}{തല കൈ (tala kai)}]{\ml{തല}, \ml{കൈ} — head and hand, matching Tamil closely}

\label{lesson:ML-C13-shareera-bhaagangal}



\begin{quote}
Malayalam's family words (\emph{achan}, \emph{cēṭṭan}...) went their own way from Tamil's. Body parts snap right back to the close match you saw in colors and days.
\end{quote}

\subsection*{What to know first}
\begin{itemize}
\raggedright
  \item \textbf{\ml{തല}} (\emph{tala}) = \textquotedblleft{}\textbf{head}\textquotedblright{} — matches Tamil \emph{talai} closely.
  \item \textbf{\ml{കൈ}} (\emph{kai}) = \textquotedblleft{}\textbf{hand}\textquotedblright{} — matches Tamil \emph{kai} almost exactly.
\end{itemize}

Both native Dravidian, both close cousins of Tamil's own words — a reminder that Malayalam and Tamil's vocabulary agreement varies concept by concept, not uniformly.

\subsection*{Your turn}
\begin{itemize}
\raggedright
  \item \emph{Say it:} \textquotedblleft{}tala\textquotedblright{} — head
  \item \emph{Say it:} \textquotedblleft{}kai\textquotedblright{} — hand
  \item \emph{Say it:} the match — nearly identical to Tamil's talai/kai
  \item \emph{Recall:} say \emph{dayavāyi}
\end{itemize}

\begin{tcolorbox}[breakable,colback=teal!4,colframe=teal!35!black,title={Before you move on}]
What are Malayalam's words for head and hand, and how closely do they match Tamil? (\textbf{\ml{തല} tala, \ml{കൈ} kai} — very closely, both native Dravidian.)
\end{tcolorbox}
\section[\texorpdfstring{\ml{ത} (ta)}{ത (ta)}]{\ml{ത} — one character, met inside words you already say}

\label{lesson:ML-S08-letter-ta}



\begin{quote}
Before the new one: \ml{മ} — what does it do?

One character this time. Just one — and you have been saying it for pages without knowing which mark on the page it was.
\end{quote}

\subsection*{Script you'll notice: \ml{ത}}
\textbf{\ml{ത}} — \emph{ta}.

It is a \textbf{consonant}, and in this script a consonant is never bare: it comes with an \emph{a} already in it. So it is not \emph{t}, it is \textbf{ta}.

You already say these, and every one of them has \ml{ത} somewhere inside it:

\begin{itemize}
\raggedright
  \item \textbf{\ml{അതെ}} — yes (athe — \textquotedblleft{}that {[}is so{]}\textquotedblright{})
  \item \textbf{\ml{എന്ത്}} — what
  \item \textbf{\ml{സന്തോഷം}} \emph{santōṣaṁ} — joy / pleased to meet you
  \item \textbf{\ml{താമസിക്കുക}} \emph{tāmasikkuka} — to live, to stay
\end{itemize}

\subsection*{Writing: \ml{ത} — follow the numbered strip}
\hlblockfigure{\detokenize{figures/ML-S08-letter-ta-filmstrip.pdf}}{How \ml{ത} is written, stroke by stroke}

Follow the numbered strip: start where movement 1 begins, and let each panel show you the next movement before your pen makes it. Then write \ml{ത} yourself — slowly, and larger than it is printed.

\begin{quote}
The strip shows \textbf{where to start the character and which way to travel}, in one attested order. That is taught with real variation from school to school, so the strip names its source beneath it: treat it as a sound way in, not the only one.
\end{quote}

\subsection*{Your turn}
\begin{itemize}
\raggedright
  \item \emph{Look:} at these words, and find \ml{ത} in the ones that have it
\end{itemize}

\begin{quote}
\ml{അതെ}  ·  \ml{എന്ത്}  ·  \ml{നമസ്കാരം}
\end{quote}

\begin{itemize}
\raggedright
  \item \emph{Trace it:} \ml{ത} three times, saying \emph{ta} as you finish each one
  \item \emph{Look:} back at any page of this chapter and find \ml{ത} once more
\end{itemize}

\begin{tcolorbox}[breakable,colback=teal!4,colframe=teal!35!black,title={Before you move on}]
Which character is this — \ml{ത}? What sound does it carry? (\emph{\textbf{ta}}.) Name one word you already say that contains it.
\end{tcolorbox}
\section[\texorpdfstring{\ml{ട} (ṭa)}{ട (ṭa)}]{\ml{ട} — one character, met inside words you already say}

\label{lesson:ML-S117-letter-tta}



\begin{quote}
Before the new one: \ml{ത} — what does it do?

One character this time. Just one — and you have been saying it for pages without knowing which mark on the page it was.
\end{quote}

\subsection*{Script you'll notice: \ml{ട}}
\textbf{\ml{ട}} — \emph{ṭa}.

It is a \textbf{consonant}, and in this script a consonant is never bare: it comes with an \emph{a} already in it. So it is not \emph{ṭ}, it is \textbf{ṭa}.

You already say these, and every one of them has \ml{ട} somewhere inside it:

\begin{itemize}
\raggedright
  \item \textbf{\ml{വീണ്ടും} \ml{കാണാം}} — we'll meet again
  \item \textbf{\ml{ഒന്ന്} \ml{രണ്ട്} \ml{മൂന്ന്} \ml{നാല്} \ml{അഞ്ച്}} — one to five
  \item \textbf{\ml{ആറ്} \ml{ഏഴ്} \ml{എട്ട്} \ml{ഒമ്പത്} \ml{പത്ത്}} — six to ten
  \item \textbf{\ml{അച്ഛൻ} \ml{അമ്മ} \ml{ചേട്ടൻ} \ml{അനിയൻ} \ml{ചേച്ചി} \ml{അനിയത്തി}} — father, mother, and four age-graded sibling words
\end{itemize}

\subsection*{Writing: \ml{ട} — follow the numbered strip}
\hlblockfigure{\detokenize{figures/ML-S117-letter-tta-filmstrip.pdf}}{How \ml{ട} is written, stroke by stroke}

Follow the numbered strip: start where movement 1 begins, and let each panel show you the next movement before your pen makes it. Then write \ml{ട} yourself — slowly, and larger than it is printed.

\begin{quote}
The strip shows \textbf{where to start the character and which way to travel}, in one attested order. That is taught with real variation from school to school, so the strip names its source beneath it: treat it as a sound way in, not the only one.
\end{quote}

\subsection*{Your turn}
\begin{itemize}
\raggedright
  \item \emph{Look:} at these words, and find \ml{ട} in the ones that have it
\end{itemize}

\begin{quote}
\ml{നമസ്കാരം}
\end{quote}

\begin{itemize}
\raggedright
  \item \emph{Trace it:} \ml{ട} three times, saying \emph{ṭa} as you finish each one
  \item \emph{Look:} back at any page of this chapter and find \ml{ട} once more
\end{itemize}

\begin{tcolorbox}[breakable,colback=teal!4,colframe=teal!35!black,title={Before you move on}]
Which character is this — \ml{ട}? What sound does it carry? (\emph{\textbf{ṭa}}.) Name one word you already say that contains it.
\end{tcolorbox}
\section[Practice]{Chapter 13 checkpoint — head, hand, \ml{ത} and \ml{ട}}

\label{lesson:ML-C13-checkpoint}



\begin{quote}
Nothing new here. The two body words out loud first, then the two letters this chapter put on the page.
\end{quote}

\subsection*{Your turn — head and hand}
\begin{itemize}
\raggedright
  \item \emph{Say it:} \textquotedblleft{}head\textquotedblright{}, then \textquotedblleft{}hand\textquotedblright{}, in Malayalam
\end{itemize}

(\emph{Tala}; \emph{kai}.)

\begin{itemize}
\raggedright
  \item How close are they to Tamil's words? (Very close: Tamil \emph{talai} and \emph{kai}. Both pairs are native Dravidian, unlike the family words, where Malayalam went its own way.)
\end{itemize}

\subsection*{Script — \ml{ത} and \ml{ട}}
\begin{itemize}
\raggedright
  \item \textbf{\ml{തല}} is \emph{tala}, \textquotedblleft{}head\textquotedblright{}. Which letter opens it, and what does it carry? (\textbf{\ml{ത}}, \emph{ta}: a consonant with its \emph{a} built in.)
  \item \textbf{\ml{രണ്ട്}} is \emph{raṇṭŭ}, \textquotedblleft{}two\textquotedblright{}. Which letter in it carries \emph{ṭa}? (\textbf{\ml{ട}}.)
\end{itemize}

\emph{Read it:} \textbf{\ml{തല}} aloud, and point to the \ml{ത} in it

\emph{Read it:} \textbf{\ml{രണ്ട്}} aloud, and point to the \ml{ട} in it

\subsection*{Writing — from sound}
\emph{Cover:} the Script block above, so no model stays in view

\emph{Write it:} the letter for \emph{ta}, then the letter for \emph{ṭa}

\emph{Check:} against the key — \textbf{\ml{ത}}, \textbf{\ml{ട}}

\emph{Write it:} one repaired shape, then stop

\begin{tcolorbox}[breakable,colback=teal!4,colframe=teal!35!black,title={Before you move on}]
Touch your head and say its word, then your hand and say its word. That is the chapter: two words Malayalam shares with Tamil, and two letters.

Next chapter: the seasons.
\end{tcolorbox}
